% Part: computability
% Chapter: computability-theory
% Section: prop-reduce

\documentclass[../../../include/open-logic-section]{subfiles}

\begin{document}

% SEGMENT OLP-0244-S01

\olfileid{cmp}{thy}{ppr}
\olsection{குறைத்தன்மையின் பண்புகள்}

$A$ ஆனது~$B$ க்குக் குறைந்தால் $A \leq_m B$ என்று எழுதுகிறோம்.
$A \leq_m B$ எனில் $A$ ஆனது~$B$ ஐவிட ``கடினமானதல்ல'' என்றும்,
$B$ ஆனது~$A$ ஐவிட ``அதே அளவு கடினமானது அல்லது மேலும்
கடினமானது'' என்றும் இக்குறியீடு உணர்த்துகிறது. மேலும், எண்கள்மீதான
வழக்கமான $\le$ போலவே, கணங்களை (அல்லது தீர்மானச் சிக்கல்களை)
அவற்றின் சிக்கற்பாட்டின்படி $\le_m$ வரிசைப்படுத்துகிறது என்றும்
அது உணர்த்துகிறது. பின்வரும் இரு முன்மொழிவுகள் இந்த உள்ளுணர்வை
ஆதரிக்கின்றன. $\le_m$ கடப்புடையது என்று முதலாவது கூறுகிறது.

\begin{prop}
\ollabel{prop:trans-red}
$A \leq_m B$ உம் $B \leq_m C$ உம் எனில், $A \leq_m C$.
\end{prop}

\begin{proof}
$A$ இலிருந்து~$B$ க்கான ஒரு குறைத்தலுடன் $B$ இலிருந்து
$C$ க்கான ஒரு குறைத்தலைச் சேர்த்தால், $A$ இலிருந்து~$C$ க்கான
ஒரு குறைத்தல் கிடைக்கிறது.
\end{proof}

\begin{prob}
$f$,~$g$ ஆகியவை முறையே $A$ இலிருந்து~$B$ க்கும், $B$ இலிருந்து
~$C$ க்கும் பலவற்றுக்கு ஒன்று குறைத்தல்கள் எனில்,
$\comp{f}{g}$ ஆனது $A$ இலிருந்து~$C$ க்கான பலவற்றுக்கு ஒன்று
குறைத்தல் என்று காட்டி \olref{prop:trans-red} ஐ நிறுவுக.
\end{prob}

\begin{prop}
\ollabel{prop:reduce}
$A$, $B$ ஆகியவை ஏதேனும் இரு கணங்களாகவும், $A \leq_m B$ ஆகவும்
இருக்கட்டும்.
\begin{enumerate}
\item $B$ கணிக்கத்தக்கவாறு எண்ணிடத்தக்கது எனில், $A$ உம் அவ்வாறே.
\item $B$ கணிக்கத்தக்கது எனில், $A$ உம் அவ்வாறே.
\end{enumerate}
\end{prop}

\begin{proof}
$f$ ஆனது $A$ இலிருந்து~$B$ க்கான பலவற்றுக்கு ஒன்று குறைத்தல்
ஆகட்டும். முதல் கூற்றுக்கு, பகுதிச் சார்பு~$g$ இன் வரையறைப்புலமாக
$B$ இருந்தால், ~$\comp{f}{g}$ இன் வரையறைப்புலமாக~$A$ இருப்பதைச்
சரிபார்த்தால் போதும்:
\begin{align*}
x \in A & \text{ எனில், எனில் மட்டுமே } f(x) \in B \\
& \text{ எனில், எனில் மட்டுமே }  g(f(x)) \fdefined.
\end{align*}

இரண்டாவது கூற்றுக்கு, $B$ கணிக்கத்தக்கது எனில் $B$ உம்
~$\Complement{B}$ உம் கணிக்கத்தக்கவாறு எண்ணிடத்தக்கவை என்பதை
நினைவுகூர்க (\olref[cmp]{thm:ce-comp}). $f$ ஆனது
$\Complement{A}$ இலிருந்து $\Complement{B}$ க்கும் பலவற்றுக்கு
ஒன்று குறைத்தல் என்பதைச் சரிபார்ப்பது கடினமல்ல. ஆகவே இந்நிரூபணத்தின்
முதல் பகுதியால் $A$ உம்~$\Complement{A}$ உம்
!!{computably enumerable}. எனவே \olref[cmp]{thm:ce-comp} ஆல்
$A$ உம் கணிக்கத்தக்கது. (மாற்றாக,
$\Char{A} = \comp{f}{\Char{B}}$ என்று சரிபார்க்கலாம்; எனவே
$\Char{B}$ கணிக்கத்தக்கது எனில்~$\Char{A}$ உம் அவ்வாறே.)
\end{proof}

\begin{prob}
$f$ ஆனது $A$ இலிருந்து~$B$ க்கான பலவற்றுக்கு ஒன்று குறைத்தல்
எனக் கொள்க. $f$ ஆனது $\Complement{A}$ இலிருந்து
$\Complement{B}$ க்கும் பலவற்றுக்கு ஒன்று குறைத்தல் என்று காட்டுக.
\end{prob}

\begin{prob}
$f\colon A \to B$ ஒரு பலவற்றுக்கு ஒன்று குறைத்தல் எனில்,
$\Char{A} = \comp{f}{\Char{B}}$ என்று காட்டுக.
\end{prob}

\begin{digress}
குறிப்பாகத் தீர்மானிக்கவியலாமை முடிவுகளை நிறுவுவதற்கு, பிற
சூழல்களில் \emph{டியூரிங் குறைத்தன்மை} எனப்படும் மேலும் பொதுவான
குறைத்தன்மைக் கருத்து பயன்படுகிறது. \olref[cmp]{cor:comp-k} ஆல்
~$K_0$ இன் நிரப்பி கணிக்கத்தக்கவாறு எண்ணிடத்தக்கது அல்ல;
எனவே அது~$K_0$ க்குக் குறைக்கத்தக்கது அல்ல என்பதைக் கவனிக்க.
ஆனால் $K_0$ பற்றிய கேள்விகளின் விடைகள் உங்களுக்குத் தெரிந்திருந்தால்,
அதன் நிரப்பி பற்றிய கேள்விகளின் விடைகளும் தெரியும் என்பது உள்ளுணர்வாகத்
தெளிவு. ~$B$ பற்றிய கேள்விகளைக் கேட்கக்கூடிய ஒரு கணிக்கத்தக்க
செயல்முறையைக் கொண்டு~$A$ பற்றிய கேள்விகளுக்கான விடைகளைத் தீர்மானிக்க
முடிந்தால், கணம் $A$ ஆனது~$B$ க்கு டியூரிங் குறைக்கத்தக்கது எனப்படும்.
பலவற்றுக்கு ஒன்று குறைத்தன்மையில் (1)~$B$ பற்றி ஒரே ஒரு கேள்வியை
மட்டுமே கேட்கலாம்; (2)~``ஆம்'' என்ற விடை $A$ பற்றிய கேள்விக்கான
``ஆம்'' என்ற விடையாகவே மாற வேண்டும், ``இல்லை'' என்பதற்கும் அவ்வாறே.
இதனுடன் ஒப்பிடுகையில் டியூரிங் குறைத்தன்மை தளர்வானது. $A$ ஆனது~$B$
க்குக் டியூரிங் குறைக்கத்தக்கதாகவும், $B$ கணிக்கத்தக்கதாகவும்
இருந்தால் $A$ உம் கணிக்கத்தக்கதாகவே இருக்கும். (ஆனால் நாம்
கண்டதுபோல், கணிக்கத்தக்கவாறு எண்ணிடத்தக்க தன்மைக்கான ஒத்த கூற்று
உண்மையல்ல.)

நாம் விவாதித்த பல்வேறு குறைத்தன்மைக் கருத்துகளைச் சிந்தித்து,
அவற்றுக்கிடையிலான வேறுபாடுகளைப் புரிந்துகொள்ள வேண்டும். எனினும்
இவ்வத்தியாயத்தில் பலவற்றுக்கு ஒன்று குறைத்தன்மையை மட்டுமே கையாள்வோம்.
கடந்த பத்தியில் விவாதிக்கப்பட்ட இருவகைக் குறைத்தன்மைகளுக்கும்
கணிப்புச் சிக்கற்பாட்டில் ஒப்பான கருத்துகள் உள்ளன; கூடுதலாக,
டியூரிங் பொறிகள் பல்லுறுப்புக்கோவை நேரத்தில் இயங்க வேண்டும்.
பலவற்றுக்கு ஒன்று குறைத்தன்மையின் சிக்கற்பாட்டுப் பதிப்பு
\emph{கார்ப் குறைத்தன்மை} என்றும், டியூரிங் குறைத்தன்மையின்
சிக்கற்பாட்டுப் பதிப்பு \emph{குக் குறைத்தன்மை} என்றும் அழைக்கப்படுகின்றன.
\end{digress}

\end{document}
